\documentclass[10pt,letterpaper]{article}

\usepackage[T1]{fontenc}
\usepackage[utf8]{inputenc}
\usepackage{microtype}
\usepackage{newpxtext, newpxmath}
\usepackage[margin=1cm]{geometry}
\usepackage{enumitem}
\usepackage{xcolor}
\usepackage[colorlinks=true,linkcolor=blue,urlcolor=blue]{hyperref}

\setlength{\parindent}{0pt}

% --- Spacing model -------------------------------------------------------
% All vertical rhythm is point-based and collapses via \addvspace so
% adjacent glue never stacks: sections get 7pt above / 2pt+rule+2pt
% below; every entry block gets 2pt above it, whether it follows text
% or a bullet list.

\newcommand{\ressection}[1]{%
  \par\addvspace{7pt}%
  {\large\bfseries #1}\par\nobreak
  \vspace{2pt}\noindent\rule{\linewidth}{1pt}\par\nobreak
  \vspace{2pt}%
}

% Bold left text, right-aligned date/location
\newcommand{\entry}[2]{%
  \par\addvspace{2pt}%
  \noindent\textbf{#1}\hfill #2\par\nobreak
}
% Plain left text (role line), right-aligned date
\newcommand{\subentry}[2]{%
  \noindent #1\hfill #2\par
}
% Single line that is only partly bold (e.g. award lines)
\newcommand{\mixedentry}[3]{%
  \par\addvspace{2pt}%
  \noindent\textbf{#1} #2\hfill #3\par
}
% Paragraph-style description under an entry (matches the website's
% Software section, which uses prose instead of bullet lists)
\newcommand{\entrydesc}[1]{%
  \par\addvspace{1pt}%
  \noindent #1\par
}

\setlist[itemize]{leftmargin=*, itemsep=0pt, topsep=1pt, parsep=0pt, partopsep=0pt}

\linespread{0.96}\selectfont

\begin{document}

% --- Header: centered stack (name / contact line / tagline), rule below ---
\begin{center}
  {\fontsize{18}{21}\selectfont\bfseries Alex Davison}\\[3pt]
  {\footnotesize Pittsburgh, PA, (919) 819--7153, \href{mailto:alex@alexdavison.com}{alex@alexdavison.com}, \href{https://www.linkedin.com/in/alexander-davison-62086127b/}{LinkedIn}, \href{https://alexdavison.com/}{alexdavison.com}}\\[3pt]
  {\small\itshape Software developer with a strong background in mathematics.}
\end{center}

\ressection{EDUCATION}

\entry{Stony Brook University}{Stony Brook, New York}
\begin{itemize}
  \item \textbf{Bachelor of Science in Mathematics,} Minor in Computer Science \hfill Expected May 2027
  \item \textbf{Courses:} Computer Proofs, Algorithms, Cryptography (Graduate), Number Theory
\end{itemize}

\ressection{SKILLS}

\begin{itemize}
  \item \textbf{Programming:} Python, TypeScript, C++
  \item \textbf{Tools:} Git, GitHub, Linux, Bash, numpy, matplotlib, \LaTeX
  \item \textbf{Math:} analysis, formal verification, algorithm design
  \item \textbf{Core Competencies:} project strategy, problem-solving, collaboration
  \item \textbf{Languages:} English (native), Mandarin (fluent), Russian (beginner)
\end{itemize}

\ressection{EXPERIENCE}

\entry{Gecko Robotics}{Summer 2026}
\subentry{Software Intern}{}
\entrydesc{Solo-built the logbook pages and integration to Gecko's Cantilever software platform for power plants, encompassing all logging modes including daily operator rounds, offloading reports, and work orders. Integrated logbook context into Cantilever's agentic document search and LLM pipeline, and interviewed control room operators on-site at multiple plants, translating operator feedback into logbook features. Rolled out to 150+ plants in the US, comprising 5\% of the North American power grid.}

\entry{Reprogramming}{}
\subentry{Software Developer (Contracted)}{2024}
\entrydesc{Designed and developed a JavaScript web application for Thales Academy, a network of 13 Pre-K--12 schools, automating creation of math worksheets for elementary students, with teacher-tailorable parameters to ensure alignment with curricula.}

\entry{DataAnnotation}{}
\subentry{Freelance Software Developer}{2023--25}
\entrydesc{Worked on a variety of script and codebase revisions across multiple languages (Python, Java, JavaScript, C, C++), providing thorough code reviews, troubleshooting approaches, and detailed explanations of revisions implemented.}

\entry{Math Tutoring}{}
\subentry{Math Tutor}{2023--now}
\entrydesc{Tutored students weekly on subjects ranging from elementary level math to calculus 4: private (2023--2024), Mathnasium (2024--2025), SBU ASTC (2025--present).}

\ressection{PROJECTS}

\entry{ChessDB Mobile}{2023}
\entrydesc{A mobile-friendly web app for exploring chess openings with live statistics and evaluation data from chessdb.cn and lichess.org. The original codebase behind the \href{https://chesspath.pro}{ChessPath} program, with hundreds of users.}

\entry{EOcross Solver}{2022}
\entrydesc{A novel algorithm for Rubik's cube EOcross-solving that balances computer optimization with human intuition. The \href{https://crystalcuber.com}{CrystalCube} project implements my algorithm, used for training by several top speedsolvers.}

\ressection{STUDENT ENGAGEMENT}

\entry{HackNYU 2025}{2025}
\begin{itemize}
  \item Put together and led a four-person team. Guided brainstorming and project planning sessions.
  \item Developed front-end user profile UI, and back-end database support in Python with SQLite.
  \item Coordinated the merging of all contributors' work in Git/GitHub and produced the finalized project.
\end{itemize}

\mixedentry{USA Computing Olympiad}{}{2017}

\mixedentry{Competitive Chess}{-- 1\textsuperscript{st} place US Junior Chess Congress 2020, Stony Brook vs.\ Yale Invitationals 2024}{2016--24}

\end{document}
